\section{通信保障下的运输与中继联合时序}
\subsection{双向链路与单跳中继}
空中中继可以补充灾后地面覆盖，但其位置、回传链路与服务时段必须同时可行\cite{zeng}。按官方通信规则，双向允许传播损耗取两个方向预算的较小值。对三维端点间距离 $D_{ij}(t)$（km）和频率 $f$（MHz），定义
\begin{align}
L_{ij}(t)&=32.45+20\log_{10}f+20\log_{10}D_{ij}(t)+L_{\rm obs}\,b_{ij}(t),\label{eq:fspl}\\
M_{ij}(t)&=L^{\max}_{i\leftrightarrow j}-L_{ij}(t).\label{eq:margin}
\end{align}
自由空间项采用与 ITU-R P.525 一致的频率、距离量纲\cite{itu}。$b_{ij}=1$ 表示视线射线与闭 DEM 像元地形相交，包括切触；有遮挡时附加题设损耗。官方参数给出 2400 MHz、遮挡附加 10 dB，直连、接入和回传预算分别为 122、116、126 dB。

令 $\mathcal R(t)$ 为已完成建链且尚未结束服务的中继集合，则运输状态可用的必要判据写为
\begin{equation}
M_p(t)=\max\left\{M_{pG}(t),\ \max_{r\in\mathcal R(t)}\min[M_{pr}(t),M_{rG}(t)]\right\}\ge0.
\label{eq:comm}
\end{equation}
实现中逐运输架次使用其配置的单个保障中继，不能把不可同时工作的不同中继的链路拼接，也不允许中继之间多跳。直连足够时可直接连接网关；需要中继时接入、回传与时间窗口必须同时成立。

\subsection{中继能量与固定离散方案的线性规划}
一个中继架次的能耗由往返巡航、正向爬升、建链与服务组成：
\begin{equation}
E_r=E_r^{\rm cruise}+E_r^{\rm climb}
 +(P_r^{\rm hover}+P_r^{\rm comm})\frac{t_r^{\rm setup}+b_r-a_r}{3600},
\qquad E_r\le(1-\rho_r)E_r^{\rm use}.
\label{eq:relay}
\end{equation}
其中两个功率以 kW 计，服务与建链时间以 s 计。中继悬停点位于 DEM 内且离地不超过 300 m。来回航段既满足地形净空，也到达选定悬停海拔；建链完成前不计为可提供服务。

配置文件固定货箱、路径、机型、资源先后关系以及中继坐标。对该离散方案，逐相位采样得到需要中继的相对时间范围 $[\alpha_p,\beta_p]$，加入
\begin{equation}
a_{r(p)}\le s_p+\alpha_p-5,\qquad
b_{r(p)}\ge s_p+\beta_p+5.
\label{eq:window}
\end{equation}
再加入式\eqref{eq:deadline}、资源顺序、建链前飞行时间、中继周转时间和式\eqref{eq:relay}。由于航段及载荷固定，上述约束对开始时刻和服务窗口均为线性。依次最小化加权延误、联合完工历时、中继服务总时长，前一层最优值以数值容差锁定后求下一层，使用 HiGHS 线性规划求解器\cite{highs}。运输能耗此时固定，中继服务时长最小化也就减少了可调整的联合能耗。

这一步只证明给定离散方案下的时序最优性。悬停点的全空间搜索、货箱组批和资源顺序的联合全局最优性没有由该 LP 证明。论文保留这一条件，防止把“LP 已最优”扩大解释为 Q3 原问题全局最优。

\subsection{联合方案与通信保障代价}
\textbf{联合候选含 \QThreeTrips{} 个运输架次与 \RelayTrips{} 个中继架次，总能耗 \JointEnergy{} kWh，联合完工 \JointTime{} s。}其中运输能耗 \QThreeTransportEnergy{} kWh，中继能耗 \RelayEnergy{} kWh；运输完工时刻为 \QThreeTransportTime{} s。加权延误为 \QThreeLate{} 优先系数·s，硬截止最小余量 \HardSlack{} s。
\paperfig{fig06_communication}{中继覆盖实际轨迹上直连不足的采样状态。橙点表示直连不足，空心三角为中继悬停点，星号为最紧通信状态；下图为 60 s 时间窗内全部在途状态的最小可用链路余量。检查包含每个飞行与交接相位的端点，共 \StatesOne{} 个状态，最小余量 \MarginOne{} dB；该结论仅针对采样集合。}{communication}

与选定 Q2 候选相比，Q3 联合能耗增加 \EnergyIncrease\%，完工历时增加 \TimeIncrease\%。这是当前两个候选在不同约束集合下的结果差异，不是通信约束的理论最小代价。图\ref{fig:communication}显示，通信困难并非只发生在目的节点；只检查起点、终点和中继位置会漏掉航段中的遮挡状态。

射线遮挡是离散地形上的非光滑判定，通信余量很小的区间尤其需要谨慎。尽管所有已检查状态通过，仍不能据此断言“任意连续时刻均满足题设”。第\ref{sec:validation}节进一步加密采样并报告余量变化。
\FloatBarrier
